\documentclass[12pt]{article}
\usepackage{amsmath, amssymb}
\usepackage{booktabs}
\usepackage{geometry}
\usepackage{enumitem}
\geometry{margin=1in}

\title{Lecture 5: Bayes’ Theorem, Random Variables, and Probability Mass Function}
\date{}
\begin{document}
\maketitle

\section{Bayes’ Theorem}

\subsection{Weighted Average of Conditional Probabilities}

Let \textbf{A} and \textbf{B} be events.

Any event \textbf{A} can be expressed as:
\[
A = AB \cup AB^c
\]

This follows because for an outcome to belong to \textbf{A}, it must either:

\begin{itemize}
\item Occur in both \textbf{A and B}, or
\item Occur in \textbf{A but not in B}
\end{itemize}

Since \textbf{AB} and \textbf{AB}$^{c}$ are mutually exclusive, by Axiom 3 of probability:
\[
Pr(A) = Pr(AB) + Pr(AB^c)
\]

Using conditional probability:
\[
Pr(A) = Pr(A|B)Pr(B) + Pr(A|B^c)[1 - Pr(B)]
\]

\subsubsection*{Interpretation:}

\begin{quote}
The probability of event \textbf{A} is a \textbf{weighted average} of conditional probabilities, where the weights are the probabilities of the conditioning events.
\end{quote}

\subsection{Learning by Example — Insurance Problem}

\subsubsection*{Example 3.1 (Part 1)}

An insurance company classifies people into:

\begin{itemize}
\item Accident-prone
\item Not accident-prone
\end{itemize}

Given:
\[
Pr(\text{Accident | accident-prone}) = 0.4
\]
\[
Pr(\text{Accident | not accident-prone}) = 0.2
\]
\[
Pr(\text{Accident-prone}) = 0.3
\]

Let:
\begin{itemize}
\item $(A_1)$ = policyholder has an accident within a year
\item $(A)$ = policyholder is accident-prone
\end{itemize}

Apply weighted average:
\[
Pr(A_1) = Pr(A_1|A)Pr(A) + Pr(A_1|A^c)Pr(A^c)
\]
\[
= (0.4)(0.3) + (0.2)(0.7)
\]
\[
= 0.12 + 0.14 = 0.26
\]

\textbf{Probability of accident = 0.26}

\subsubsection*{Example 3.1 (Part 2)}

Given that the policyholder \textbf{had an accident}, find the probability they are accident-prone:
\[
Pr(A|A_1)
\]

Using Bayes’ formula:
\[
Pr(A|A_1) = \frac{Pr(A)Pr(A_1|A)}{Pr(A_1)}
\]
\[
= \frac{(0.3)(0.4)}{0.26} = \frac{6}{13}
\]

\subsection{Law of Total Probability}

Let:
\[
B_1, B_2, ..., B_n
\]

be \textbf{mutually exclusive events} such that:
\[
\bigcup_{i=1}^n B_i = B
\]

Since:
\[
A = \bigcup_{i=1}^n AB_i
\]

and the events $(AB_i)$ are mutually exclusive:
\[
Pr(A) = \sum_{i=1}^n Pr(AB_i)
\]

Using conditional probability:
\[
Pr(A) = \sum_{i=1}^n Pr(A|B_i)Pr(B_i)
\]

This is called the \textbf{Law of Total Probability (Formula 3.4)}

\subsection{Bayes’ Formula (Proposition 3.1)}

Using:
\[
Pr(AB_i) = Pr(B_i|A)Pr(A)
\]

we obtain:
\[
Pr(B_i|A) = \frac{Pr(A|B_i)Pr(B_i)}{\sum_{j=1}^n Pr(A|B_j)Pr(B_j)}
\]

\subsubsection*{Terminology:}

\begin{itemize}
\item $Pr(B_i)$ $\rightarrow$ \textit{a priori probability}
\item $Pr(B_i | A)$ $\rightarrow$ \textit{posterior probability}
\end{itemize}

\subsection{Bayes Example — Colored Cards}

Three cards:

\begin{center}
\begin{tabular}{ll}
\toprule
Card & Sides \\
\midrule
1 & Red – Red \\
2 & Black – Black \\
3 & Red – Black \\
\bottomrule
\end{tabular}
\end{center}

One card selected randomly.

Let:
\begin{itemize}
\item RR = all red card
\item BB = all black card
\item RB = mixed card
\item R = upper side is red
\end{itemize}

We seek:
\[
Pr(RB | R)
\]

Using Bayes:
\[
Pr(RB|R)=\frac{Pr(R|RB)Pr(RB)}{Pr(R|RR)Pr(RR)+Pr(R|RB)Pr(RB)+Pr(R|BB)Pr(BB)}
\]

Substitute:
\[
=\frac{(1/2)(1/3)}{(1)(1/3)+(1/2)(1/3)+(0)(1/3)}
\]
\[
=\frac{1/6}{1/3+1/6}=\frac{1}{3}
\]

\textbf{Correct probability = 1/3}

Common mistake: assuming two cases are equally likely (which they are not).

\section{Random Variables}

\subsection{Motivation}

Often we are interested not in outcomes themselves but in a \textbf{numerical function of outcomes}.

Examples:
\begin{itemize}
\item Sum of dice rather than each die
\item Number of heads rather than full coin sequence
\end{itemize}

\subsection{Definition}

A \textbf{random variable} is a real-valued function defined on the sample space:
\[
X: \Omega \rightarrow \mathbb{R}
\]

which assigns a real number to each sample point.

The lecture restricts attention to \textbf{discrete random variables}:

\begin{itemize}
\item Values are finite or countably infinite
\item Actual values form a discrete subset of $\mathbb{R}$
\end{itemize}

\subsection{Distribution of a Random Variable}

Two components:

\begin{enumerate}
\item Values the variable can take
\item Probabilities of those values
\end{enumerate}

For any value $(a)$:
\[
\{ \omega \in \Omega : X(\omega)=a \}
\]

is an event, abbreviated as:
\[
X = a
\]

The collection:
\[
Pr(X=a)
\]

for all values of $(a)$ is called the \textbf{distribution of X}.

\subsection{Visualization}

The distribution can be represented by a \textbf{bar diagram}:

\begin{itemize}
\item x-axis $\rightarrow$ possible values
\item bar height $\rightarrow$ $(Pr[X=a])$
\end{itemize}

\subsection{Example — Three Coin Tosses}

Let:
\[
Y = \text{number of heads}
\]

Possible values:
\[
0,1,2,3
\]

Probabilities:
\[
P(Y=0)=\frac{1}{8}
\]
\[
P(Y=1)=\frac{3}{8}
\]
\[
P(Y=2)=\frac{3}{8}
\]
\[
P(Y=3)=\frac{1}{8}
\]

Since Y must take one of these values:
\[
1 = \sum_{i=0}^3 P(Y=i)
\]

\section{Probability Mass Function (PMF)}

\subsection{Discrete Random Variable}

A random variable that takes \textbf{at most countably many values} is called \textbf{discrete}.

\subsection{Definition of PMF}

Let:
\[
R_X = \{x_1, x_2, x_3, ...\}
\]

be the range of X.

The function:
\[
p(x_k) = Pr(X = x_k)
\]

is called the \textbf{Probability Mass Function (PMF)}.

Since X must take one of its values:
\[
\sum_k p(x_k) = 1
\]

\subsection{PMF Example (Poisson-type Form)}

Given:
\[
p(i)=c\frac{\lambda^i}{i!}, \quad i=0,1,2,...
\]

First use normalization:
\[
\sum_{i=0}^{\infty} p(i)=1
\]
\[
c \sum_{i=0}^{\infty} \frac{\lambda^i}{i!}=1
\]

Using:
\[
e^\lambda = \sum_{i=0}^{\infty} \frac{\lambda^i}{i!}
\]

So:
\[
c e^\lambda = 1 \Rightarrow c = e^{-\lambda}
\]

\subsubsection*{Find $(P(X=0))$}

\[
P(X=0)=e^{-\lambda}\frac{\lambda^0}{0!}=e^{-\lambda}
\]

\subsubsection*{Find $(P(X>2))$}

\[
P(X>2)=1-P(X\le2)
\]
\[
=1-[P(X=0)+P(X=1)+P(X=2)]
\]

\end{document}


\end{document}
